\section{Desenho experimental}
\label{sec:experimentos}

\subsection{Objetivo e hipóteses operacionais}
Ilustrar \textbf{análise de sensibilidade} sob um modelo transparente: mantemos fixos o backlog inicial, a semente do PRNG, a política de WIP e os parâmetros globais (Tabela~\ref{tab:param}), variando apenas (i) sinergia em um par colaborativo em \texttt{dev}, (ii) sinergia negativa no handoff entre analista e desenvolvedor principal, e (iii) um defeito individual (\emph{Hiperfoco}) no desenvolvedor principal.

As hipóteses \emph{operacionais} associadas a este desenho são modestas e \textbf{dependentes dos parâmetros}: esperamos que (H1) sinergia positiva no par colaborativo reduza, em média, o tempo de ciclo dos cartões que se beneficiam de maior produtividade efetiva em \texttt{dev}; (H2) sinergia negativa no handoff da analista para o desenvolvedor principal aumente atrasos via combinação de multiplicador de handoff e maior probabilidade de retrabalho; (H3) o defeito \emph{Hiperfoco} reduza throughput médio ao restringir o suporte do dado diário, com impacto variável na média de ciclo conforme a realização estocástica e a contenção do sistema. A Seção~\ref{sec:resultados} confronta essas expectativas com números gerados automaticamente.

\subsection{Equipe e backlog}
Utilizamos quatro membros: uma analista, dois desenvolvedores (para permitir cartões colaborativos) e um testador. O backlog possui sete cartões com pontuações heterogêneas; dois cartões são colaborativos com assignees nos dois desenvolvedores. Essa composição força, de forma controlada, o uso do mecanismo de colaboração em \texttt{dev} sem alterar o tamanho total do trabalho entre cenários.

\subsection{Cenários}
\begin{itemize}
  \item \textbf{A (linha de base):} sinergias nulas e sem traços.
  \item \textbf{B (colaboração):} $s=0{,}88$ entre os dois desenvolvedores; demais sinergias nulas.
  \item \textbf{C (handoff):} $s=-0{,}82$ entre analista e desenvolvedor principal; demais sinergias nulas.
  \item \textbf{D (traço):} como A, porém o desenvolvedor principal recebe o defeito \emph{Hiperfoco}.
\end{itemize}

\subsection{Parâmetros numéricos fixos}
\begin{table}[htbp]
  \centering
  \caption{Parâmetros globais usados nos experimentos (gerados a partir do código).}
  \label{tab:param}
  % Gerado por scripts/generate-paper-figures.ts — não editar à mão
\begin{tabular}{@{}lr@{}}
  \toprule
  Parâmetro & Valor \\
  \midrule
  \texttt{daysPerSprint} & 10 \\
  \texttt{numSprints} & 10 \\
  \texttt{seed} & 424242 \\
  \texttt{wipPerColumn} & 3 \\
  \texttt{planningPullMax} & 6 \\
  \texttt{synergyBeta} & 0.38 \\
  \texttt{synergyGamma} & 0.42 \\
  \texttt{collabEffMin} / \texttt{collabEffMax} & 0.72 / 1.38 \\
  \texttt{handoffEffMin} / \texttt{handoffEffMax} & 0.68 / 1.28 \\
  \texttt{handoffReworkSynergyThreshold} & -0.12 \\
  \texttt{reworkUnits} & 2 \\
  \bottomrule
\end{tabular}

\end{table}

\subsection{Métricas reportadas e reprodutibilidade}
Para cada cenário, registramos: (i) série diária de contagens por coluna (base para visualizações estilo CFD); (ii) tempos de ciclo por cartão concluído (de \texttt{ready} a \texttt{deploy}), com média e mediana por cenário; (iii) throughput por sprint, atribuindo cada conclusão ao sprint do dia global correspondente. Todos os arquivos em \texttt{paper/data/} e \texttt{paper/figures/} podem ser regenerados a partir do mesmo commit do código, o que reduz risco de descompasso entre texto e implementação.
